% Part: sets-functions-relations
% Chapter: functions
% Section: inverses

\documentclass[../../../include/open-logic-section]{subfiles}

\begin{document}

\olfileid{sfr}{fun}{inv}
\olsection{ప్రమేయాల విలోమాలు}

\begin{explain}
ప్రమేయాలను ప్రతిచిత్రణాలుగా భావిస్తాం. కాబట్టి ఆ ప్రతిచిత్రణాన్ని
``వెనక్కి తిప్పగలమా'' అనేది సహజంగా తలెత్తే ప్రశ్న.
ఉదాహరణకు, ఉత్తరవర్తి ప్రమేయం $f(x) = x + 1$ను వెనక్కి తిప్పవచ్చు:
$g(y) = y - 1$ అనే ప్రమేయం, $f$ చేసినదాన్ని ``రద్దు చేస్తుంది''.

అయితే జాగ్రత్త అవసరం. $g$ నిర్వచనం ఒక $\Int \to \Int$ ప్రమేయాన్ని
నిర్వచించినా, అది $\Nat \to \Nat$ అనే \emph{ప్రమేయాన్ని} నిర్వచించదు:
$g(0) \notin \Nat$ కదా. కనుక సరళమైన సందర్భాల్లో కూడా, ప్రమేయాన్ని
వెనక్కి తిప్పగలమో లేదో వెంటనే స్పష్టం కాదు; అది ప్రవేశం,
సహప్రవేశంపై ఆధారపడవచ్చు.

ప్రమేయపు విలోమం అనే భావన దీన్ని మరింత కచ్చితంగా వివరిస్తుంది.
\end{explain}

\begin{defn}
ప్రతి $x \in A$కూ, $y \in B$కూ $f(g(y)) = y$, $g(f(x)) = x$
అయితే, $g \colon B \to A$ అనే ప్రమేయాన్ని
$f \colon A \to B$ యొక్క \emph{విలోమం} అంటాం.
\end{defn}

$f$కు $g$ అనే విలోమం ఉంటే, తరచూ $g$కు బదులు $f^{-1}$ అని రాస్తాం.

\begin{explain}
ఏ సందర్భాల్లో ప్రమేయాలకు విలోమాలు ఉంటాయో ఇప్పుడు నిర్ణయిద్దాం.
$f\colon A \to B$ యొక్క విలోమంగా పరిగణించడానికి సహజంగా తోచేది,
కింది విధంగా ``నిర్వచించిన'' $g\colon B \to A$:
\[
g(y) = \text{$f(x) = y$ అయ్యే ``ఆ ఏకైక'' $x$.}
\]
కానీ ``నిర్వచించిన'', ``ఆ ఏకైక'' అన్న వాటికి ఉల్లేఖన చిహ్నాలు
పెట్టడమే, ఇది నిజంగా నిర్వచనం కాదని సూచిస్తుంది.
కనీసం ఇది పూర్తి సాధారణతతో ఎల్లప్పుడూ పనిచేయదు.
ఈ నిర్వచనం ఒక ప్రమేయాన్ని పేర్కొనాలంటే, $f(x) = y$ అయ్యే
ఒకే ఒక్క $x$ ఉండాలి: $g$ ఇచ్చే ఫలితం ఏకైకంగా నిర్ణయించబడాలి.
అంతేకాక, ప్రతి $y \in B$కూ అది నిర్ణయించబడాలి.
$x_1$, $x_2 \in A$ ఉండి, $x_1 \neq x_2$ అయినా $f(x_1) = f(x_2)$ అయితే,
$y = f(x_1) = f(x_2)$ వద్ద $g(y)$ ఏకైకంగా నిర్ణయించబడదు.
$f(x) = y$ అయ్యే $x$ అసలే లేకపోతే, $g(y)$ అసలే నిర్ణయించబడదు.
మరో మాటలో, $g$ నిర్వచితమవ్వాలంటే $f$, !!{injective}, !!{surjective} రెండూ కావాలి.

దీన్ని నెమ్మదిగా పరిశీలిద్దాం. ప్రశ్నను రెండుగా విడదీద్దాం:
$f\colon A \to B$ అనే ప్రమేయం ఇచ్చినప్పుడు, $g(f(x)) = x$ అయ్యే
$g\colon B \to A$ అనే ప్రమేయం ఎప్పుడు ఉంటుంది?
ఇటువంటి $g$, $f$ చేసినదాన్ని ``రద్దు చేస్తుంది'';
దీన్ని $f$ యొక్క \emph{ఎడమ విలోమం} అంటాం.
రెండో ప్రశ్న: $f(h(y)) = y$ అయ్యే $h\colon B \to A$ అనే ప్రమేయం
ఎప్పుడు ఉంటుంది? ఇటువంటి $h$ను $f$ యొక్క \emph{కుడి విలోమం} అంటాం:
$h$ చేసినదాన్ని $f$ ``రద్దు చేస్తుంది''.
\end{explain}

\begin{prop}
$f\colon A \to B$ అనేది !!{injective} అయితే, ప్రతి $x \in A$కూ
$g(f(x)) = x$ అయ్యే $f$ యొక్క \emph{ఎడమ విలోమం} $g\colon B \to A$ ఉంటుంది.
\end{prop}

\begin{proof}
$f\colon A \to B$ అనేది !!{injective} అనుకుందాం.
ఒక $y \in B$ను పరిగణించండి. $y \in \ran{f}$ అయితే,
$f(x) = y$ అయ్యే $x \in A$ ఉంటుంది.
$f$ అనేది !!{injective} కనుక, అటువంటి $x \in A$ ఒకటే ఉంటుంది.
అప్పుడు $g(y) = x$ అని నిర్వచించవచ్చు; అంటే $g(y)$ అనేది
$f(x) = y$ అయ్యే ``ఆ ఏకైక'' $x \in A$.
$y \notin \ran{f}$ అయితే, దాన్ని ఏదైనా $a \in A$కు ప్రతిచిత్రించవచ్చు.
కాబట్టి ఒక $a \in A$ను ఎంచుకుని, $g\colon B \to A$ను ఇలా నిర్వచించవచ్చు:
\[
g(y) = \begin{cases}
    x & \text{$f(x) = y$ అయితే}\\
    a & \text{$y \notin \ran{f}$ అయితే.}
\end{cases}
\]
ఇది ప్రతి $y \in B$కూ నిర్వచితమవుతుంది: అటువంటి ప్రతి $y \in \ran{f}$కూ
$f(x) = y$ అయ్యే కచ్చితంగా ఒక $x \in A$ ఉంటుంది.
నిర్వచనం ప్రకారం, $y = f(x)$ అయితే $g(y) = x$;
అంటే $g(f(x)) = x$.
\end{proof}

\begin{prob}
$f\colon A \to B$కు $g$ అనే ఎడమ విలోమం ఉంటే,
$f$ అనేది !!{injective} అని చూపండి.
\end{prob}

\begin{prop}
$f\colon A \to B$ అనేది !!{surjective} అయితే, ప్రతి $y \in B$కూ
$f(h(y)) = y$ అయ్యే $f$ యొక్క \emph{కుడి విలోమం} $h\colon B \to A$ ఉంటుంది.
\end{prop}

\begin{proof}
$f\colon A \to B$ అనేది !!{surjective} అనుకుందాం.
ఒక $y \in B$ను పరిగణించండి. $f$ అనేది !!{surjective} కనుక,
$f(x_y) = y$ అయ్యే $x_y \in A$ ఉంటుంది.
అప్పుడు $h(y) = x_y$ అని నిర్వచించవచ్చు: ప్రతి $y \in B$కూ
$f(x) = y$ అయ్యే ఏదో ఒక $x \in A$ను ఎంచుకుంటాం.
$f$ అనేది !!{surjective} కనుక, ఎంచుకోవడానికి కనీసం ఒకటి ఎల్లప్పుడూ
ఉంటుంది.\footnote{$f$ అనేది !!{surjective} కనుక, ప్రతి $y \in B$కూ
$\Setabs{x}{f(x) = y}$ అనే సమితి శూన్యేతరం.
$h$ను నిర్వచించడానికి ఈ సమితులలో ప్రతిదాని నుంచీ ఒకే $x$ను
ఎంచుకోవాలి. ఇది ఎల్లప్పుడూ సాధ్యమన్నది నిజానికి స్పష్టం కాదు:
ఈ ఎంపికలు చేయగలమన్న విషయాన్ని ఒక స్వీకృతంగా భావిస్తున్నాం.
అంటే ఈ ప్రతిపాదన ఎంపిక స్వీకృతం అని పిలిచే సూత్రాన్ని ఆధారంగా తీసుకుంటుంది;
ఈ విషయాన్ని \oliflabeldef{sth:choice::chap}{\olref[sth][choice][]{chap}లో
తిరిగి పరిశీలిస్తాం}{ఇక్కడ వివరించకుండా దాటవేస్తాం}.
అయితే అనేక ప్రత్యేక సందర్భాల్లో, ఉదాహరణకు $A = \Nat$ అయినప్పుడు,
లేదా అది పరిమితమైనప్పుడు, లేదా $f$ అనేది !!{bijective} అయినప్పుడు,
ఎంపిక స్వీకృతం అవసరం లేదు. (ముఖ్యంగా $f$ అనేది !!{bijective} అయితే,
ప్రతి $y \in B$కూ $\Setabs{x}{f(x) = y}$ అనే సమితిలో కచ్చితంగా ఒక
!!{element} ఉంటుంది; కాబట్టి ఎంచుకోవాల్సిన అవసరమే ఉండదు.)}
నిర్వచనం ప్రకారం $x = h(y)$ అయితే $f(x) = y$;
అంటే ప్రతి $y \in B$కూ $f(h(y)) = y$.
\end{proof}

\begin{prob}
$f\colon A \to B$కు $h$ అనే కుడి విలోమం ఉంటే,
$f$ అనేది !!{surjective} అని చూపండి.
\end{prob}

\begin{explain}
ముందటి నిరూపణలోని భావాలను కలిపితే, ప్రతి !!{bijection}కూ ఒక విలోమం
ఉంటుందని తెలుస్తుంది: అంటే $f$కు ఎడమ విలోమమూ కుడి విలోమమూ
రెండూ అయ్యే ఒకే ప్రమేయం ఉంటుంది.
\end{explain}

\begin{prop}\ollabel{prop:bijection-inverse}
$f\colon A \to B$ అనేది !!{bijective} అయితే,
$f^{-1}\colon B \to A$ అనే ప్రమేయం ఉంటుంది:
ప్రతి $x \in A$కూ $f^{-1}(f(x)) = x$ అనీ,
ప్రతి $y \in B$కూ $f(f^{-1}(y)) = y$ అనీ నెరవేరుతాయి.
\end{prop}

\begin{proof}
అభ్యాసం.
\end{proof}

\begin{prob}
\olref[sfr][fun][inv]{prop:bijection-inverse}ను నిరూపించండి.
మీరు $f^{-1}$ను నిర్వచించి, అది ఒక ప్రమేయమని చూపి,
అది $f$కు విలోమమని చూపాలి: అంటే ప్రతి $x \in A$కూ, $y \in B$కూ
$f^{-1}(f(x)) = x$, $f(f^{-1}(y)) = y$ అని చూపాలి.
\end{prob}

\begin{explain}
విలోమాలను పొందడానికి కొంచెం సాధారణమైన మరో మార్గం ఉంది.
ప్రతి $x \in A$కూ $f'(x) = f(x)$ అని తీసుకుంటే, ప్రతి ప్రమేయం $f$
!!a{surjection} $f' \colon A \to \ran{f}$ను ఉత్పన్నం చేస్తుందని
\olref[kin]{sec}లో చూశాం. $f$ అనేది !!{injective} అయితే,
$f'$ అనేది !!{bijective} అవుతుందనేది స్పష్టమే.
కాబట్టి \olref{prop:bijection-inverse} ప్రకారం దానికి ఏకైక విలోమం ఉంటుంది.
సంకేతాన్ని స్వల్పంగా సడలించి, కొన్నిసార్లు $f'$ విలోమాన్ని
కేవలం ``$f$ యొక్క విలోమం'' అని పిలుస్తాం.
\end{explain}

\begin{prop}\ollabel{prop:left-right}%
$f\colon A \to B$కు $g$ అనే ఎడమ విలోమమూ $h$ అనే కుడి విలోమమూ
ఉంటే, $h = g$ అని చూపండి.
\end{prop}

\begin{proof}
అభ్యాసం.
\end{proof}

\begin{prob}
\olref[sfr][fun][inv]{prop:left-right}ను నిరూపించండి.
\end{prob}

\begin{prop}\ollabel{prop:inverse-unique}
ప్రతి ప్రమేయం $f$కు గరిష్ఠంగా ఒక విలోమమే ఉంటుంది.
\end{prop}

\begin{proof}
$g$, $h$ రెండూ $f$కు విలోమాలు అనుకుందాం.
అప్పుడు ముఖ్యంగా $g$, $f$కు ఎడమ విలోమం; $h$ కుడి విలోమం.
\olref{prop:left-right} ప్రకారం $g = h$.
\end{proof}

\end{document}
